\section{Threat model, limitations, and interpretation}
The theory and the evidence have different scopes.  Proposition~1 and its
factorization corollary hold for any finite exact-terminal decision problem, but
they identify what information and interface make a decision attainable; they
do not show that a deployed verifier possesses the target-sufficient
representation.  The empirical layers establish only the representations,
relations, and case distributions that were frozen for those layers.

First, the V2/V3 campaigns use 420 synthetic mechanical-gold cases; they do not establish performance on naturally occurring scientific disagreements. The written human-adjudication rubric was never triggered because all final cases were mechanically decidable, so no claim is made about adjudicator agreement. P4-X adds five heterogeneous artifact domains and eight archetypes but remains a 400-case exact-contract study, not a population sample of scientific disputes. Second, scoring is deterministic and uses no LLM provider. It isolates authority-transition semantics but does not measure the error distribution of a full stochastic research agent.

Third, comparator mechanisms are matched reimplementations. They instantiate published ideas under one packet and budget, but the result must not be read as a leaderboard against the original ProvenAI, ProvenanceGuard, DeepSciVerify, FIRE, AttributionBench, CLAIM-BENCH, RewardHackingAgents, or related implementations. Fourth, the clean controls are intentionally unambiguous; 60/60 coverage shows that the governed pipeline did not win by blanket refusal but does not estimate recall on difficult benign cases. Fifth, the hostile taxonomy is finite. Adaptive attacks, compromised host infrastructure, side channels outside retained syscall/network telemetry, and novel evaluator attacks remain outside this campaign.

Sixth, and most consequential for how the three hypotheses should be read, two of them are measured on axes that this battery
saturates, and for H3 the mechanism is identified rather than suspected. Clean coverage is $1.0$ for all eleven panel systems and for all eight ablations,
and the correct-undetermined rate is $1.0$ for all eleven systems and is not instrumented in the
ablations at all. Only false authority promotion varies: it takes six distinct values across
the panel and four across the ablations. The whole empirical separation reported here rests on
that one axis. H1 is therefore a measurement; H2's non-inferiority pass and H3's null are
statements about a battery in which neither quantity was ever observed to move. Testing them
would require clean positives that a weaker governance mechanism actually declines, and
insufficient-evidence cases that a weaker mechanism actually over-promotes or over-abstains on
--- cases the current gold construction does not contain. On this battery, therefore, H2 and H3
should be read as design limits rather than as comparative findings.

For H3 such cases now exist, and Section~\ref{sec:results} reports them: on the repaired battery
the correct-undetermined rate takes three values across the same eleven systems, so the axis moves
and the hypothesis is decided by the systems rather than by the construction. Three things that
result does not do. It does not retrospectively license the null above --- the repaired battery is a
different battery, and the numbers reported for the frozen campaign stand as the record of what
that construction produced. It does not widen what H3 claims: nine of the ten comparator mechanisms
score zero because they cannot express an undetermined terminal at all, so the quantity measured is
terminal expressiveness under a non-compensatory gate lattice, and the margin is $1.0$ against the
strongest comparator but $0.5$ against the one mechanism that can escalate. And it does nothing for
H2, whose clean coverage is $1.0$ for every system on both batteries; that axis remains a design
limit with no case yet built that could move it.

The V3 nuisance audit is also claim-axis specific.  Its fourteen frozen probes
and thirteen seeds clear the exact \textsc{CannotCheck} axis at informedness
$0.0$, while the adjudication receipt retains four off-axis residuals and
explicitly withholds whole-register clearance.  More importantly, nine of ten
comparator mechanisms cannot emit \textsc{CannotCheck}.  In the terminology of
Section~\ref{sec:axis-theory}, V3 has nonzero observed resolution but unmatched
terminal attainability.  It therefore repairs the V2 construction failure
without yet supporting a general scientific-judgement claim.

Seventh, the mechanisms composed here are not claimed as novel, and the boundary is worth stating once and plainly. Provenance tracking, attribution evaluation, citation checking, iterative verification, evidence escalation, contamination detection, evaluator-tamper detection, abstention, benchmark auditing, claim-level auditability, provenance ontologies, and generic behavioral-integrity checking are all prior work, cited in Section~\ref{sec:related}. The relevance-versus-authorized-evidence distinction is likewise established there. What is measured here is the behavior of the non-compensatory conjunction of those prerequisites under a protected, gold-blind battery, and nothing beyond it.

P4-X does not remove that boundary.  It grants the strongest comparison product
ordinary provenance, verification, exact artifact/version binding, evaluator
custody, epoch handling, and generic authorization before testing the target
scientific-promotion relation.  Its 400/400 versus 250/400 result shows that the
registered generic relation is insufficient on the exact successor contracts;
the 400/400 ideal-product tie simultaneously shows that an information-equivalent
donor product is sufficient.  Neither count licenses a centralization advantage,
universal necessity of the exact gate list, or superiority over deployed donor
systems.

Recent benchmark-auditing work is an additional caution: expert-built tasks and evaluation harnesses can contain specification, environment, ground-truth, or grading defects that materially alter measured rankings \citep{benchguard2026,autobenchmarkaudit2026,hal2025}. Deterministic generation, typed-panel validation, exact split/harness binding, protected custody, post-run public slicing, and independent reproduction mitigate this risk, but these controls do not prove that the 13-family taxonomy is complete or that every benchmark-design assumption is correct.

The frozen naturalistic successor, identified as
\idt{P4.H3.NATURALISTIC.V1}, is the upward test rather than a claimed result.  It
specifies 768 independent source--case-family clusters in 32 strata (four
domains by eight unidentifiability mechanisms), pairs every unidentifiable unit
with an identifiable same-source control, freezes four
\textsc{CannotCheck}-capable primary comparators and a worst-domain gate, and
requires independent evaluator/raw-case custody plus a source-disjoint
replication panel.  Its outcomes have not been accessed and the external panel
is missing.  No naturalistic, cross-domain, external-comparator, or source-
disjoint result is inferred from the protocol or its power calculation.

The appropriate current interpretation is a verification-axis envelope.  V2
measures false-promotion separation and records a saturated negative axis; V3
shows that a shortcut-resistant exact axis can still be interface-limited; and
P4-X shows both failure of a registered generic-authorization relation and exact
equivalence when a donor product receives target-sufficient state and the same
predicate.  Together they support a fail-closed scientific-promotion interface
and the factorization theory that explains its limits.  They do not prove
universal scientific truth, naturalistic scientific judgement, general
abstention ability, or evaluator security.
